As \(\Tf\) increases from \(1\) to \(256\), the number of exact mapping calls per commodity decreases by \(98.9\%\) on Milan and \(98.5\%\) on the Random topology, while the amortized mapping computation time per slot decreases by \(98.8\%\) and \(98.1\%\), respectively. Since any fixed finite \(\Tf\) preserves throughput optimality, increasing the frame length does not change the supported throughput region. In the finite horizon simulations, the measured \(\hat\lambda\) differs by at most one step of the tested arrival rate grid across the frame lengths considered; these small differences are primarily finite horizon effects. For the Random topology, the \(0.55\) at \(\Tf=1\) here comes from the \(10^5\)-slot frame length configuration, whereas the \(0.52\) in Section~\ref{subsubsec:queue_blind_mapping} comes from the wireless mapping configuration with a simulation horizon of \(2\times10^4\) slots; the two experiments use different rate grids. Completion delay increases with \(\Tf\), particularly at high offered rates. Thus, frame reuse reduces mapping computation at the cost of delay rather than throughput.
